% ═══════════════════════════════════════════════════════════════
% MUSTERLÖSUNG: Describir fotos
% Klasse 8 · Spanisch · 26 BE
% ═══════════════════════════════════════════════════════════════

\documentclass[10pt, a4paper]{article}

\usepackage[utf8]{inputenc}
\usepackage[T1]{fontenc}
\usepackage[ngerman]{babel}

\usepackage{palatino}
\usepackage{helvet}
\newcommand{\sanstext}[1]{{\fontfamily{phv}\selectfont #1}}

\usepackage{microtype}
\usepackage[margin=1.4cm, top=1.6cm, bottom=1.3cm]{geometry}
\usepackage{enumitem}
\usepackage{tabularx}
\usepackage{booktabs}
\usepackage{array}
\usepackage{multicol}
\usepackage{tikz}
\usetikzlibrary{calc,positioning}

\usepackage{fancyhdr}
\usepackage{lastpage}
\usepackage{needspace}

\usepackage{xcolor}
\usepackage{tcolorbox}
\tcbuselibrary{skins}
\usepackage{amssymb}

% ═══════════════════════════════════════════════════════════════
% FARBEN
% ═══════════════════════════════════════════════════════════════

\definecolor{kopfzeile}{HTML}{1A1A1A}
\definecolor{akzent}{HTML}{C62828}
\definecolor{hellgrau}{HTML}{F2F2F2}
\definecolor{mittelgrau}{HTML}{777777}
\definecolor{dunkelgrau}{HTML}{444444}
\definecolor{loesung}{HTML}{C62828}

% ═══════════════════════════════════════════════════════════════
% VARIABLEN
% ═══════════════════════════════════════════════════════════════

\newcommand{\thema}{Describir fotos}
\newcommand{\klasse}{8}
\newcommand{\maxpunkte}{26}

% ═══════════════════════════════════════════════════════════════
% KOPF- UND FUSSZEILE
% ═══════════════════════════════════════════════════════════════

\pagestyle{fancy}
\fancyhf{}
\renewcommand{\headrulewidth}{0pt}
\renewcommand{\footrulewidth}{0pt}
\cfoot{\sanstext{\footnotesize\textcolor{mittelgrau}{\thepage\,/\,\pageref{LastPage}}}}

% ═══════════════════════════════════════════════════════════════
% BOXEN
% ═══════════════════════════════════════════════════════════════

\newtcolorbox{grammatikbox}[1][]{
    colback=hellgrau,
    colframe=hellgrau,
    boxrule=0pt,
    arc=0pt,
    left=8pt, right=8pt, top=8pt, bottom=6pt,
    borderline north={1.5pt}{0pt}{dunkelgrau},
    title={\sanstext{\small\bfseries #1}},
    coltitle=dunkelgrau,
    attach boxed title to top left={yshift=-2pt, xshift=6pt},
    boxed title style={colback=hellgrau, colframe=hellgrau, boxrule=0pt}
}

\newtcolorbox{vokabelbox}[1][]{
    colback=white,
    colframe=mittelgrau,
    boxrule=0.5pt,
    arc=0pt,
    left=6pt, right=6pt, top=6pt, bottom=4pt,
    title={\sanstext{\small\bfseries #1}},
    coltitle=dunkelgrau,
    attach boxed title to top left={yshift=-2pt, xshift=6pt},
    boxed title style={colback=white, colframe=white, boxrule=0pt}
}

\newtcolorbox{checklistbox}{
    colback=white,
    colframe=akzent,
    boxrule=1pt,
    arc=0pt,
    left=8pt, right=8pt, top=6pt, bottom=6pt
}

% ═══════════════════════════════════════════════════════════════
% BEFEHLE
% ═══════════════════════════════════════════════════════════════

\newcommand{\es}[1]{\textbf{#1}}
\newcommand{\de}[1]{{\small\textit{\textcolor{mittelgrau}{#1}}}}
\newcommand{\stamm}[1]{\textcolor{akzent}{#1}}
\newcommand{\luecke}[1]{\rule{#1}{0.4pt}}
\newcommand{\punktebox}[1]{\hfill\sanstext{\small[\_\_/#1]}}

% LÖSUNG in Rot
\newcommand{\lsg}[1]{\textcolor{loesung}{\textbf{#1}}}

\newcommand{\aufgabe}[2]{%
    \needspace{4\baselineskip}%
    \vspace{10pt}%
    \noindent\sanstext{\textbf{#1}\quad#2}%
    \vspace{5pt}%
    \nopagebreak[4]%
}

\newlist{teil}{enumerate}{1}
\setlist[teil]{label=\alph*), leftmargin=1.8em, itemsep=2pt, parsep=0pt, topsep=2pt}

\setlength{\parindent}{0pt}
\setlength{\parskip}{3pt}
\setlength{\columnsep}{20pt}

% ═══════════════════════════════════════════════════════════════
% DOKUMENT
% ═══════════════════════════════════════════════════════════════

\begin{document}

% ═══════════════════════════════════════════════════════════════
% HEADER
% ═══════════════════════════════════════════════════════════════

\noindent
\begin{tikzpicture}[remember picture, overlay]
    \fill[kopfzeile] (current page.north west) rectangle +(\paperwidth, -1cm);
    \node[anchor=west, text=white, font=\sanstext\large\bfseries]
        at ([xshift=1.4cm, yshift=-0.5cm]current page.north west) {\thema{} -- \textcolor{loesung}{MUSTERLÖSUNG}};
    \node[anchor=east, text=white, font=\sanstext\small]
        at ([xshift=-1.4cm, yshift=-0.5cm]current page.north east)
        {Spanisch\,·\,Klasse \klasse\,·\,Arbeitsblatt};
\end{tikzpicture}

\vspace{0.5cm}

\noindent
\sanstext{\small\textbf{Name:}} \luecke{5cm} \hfill
\sanstext{\small\textbf{Datum:}} \luecke{2.5cm}

\vspace{6pt}

% ═══════════════════════════════════════════════════════════════
% G1: GRAMMATIK – Presente continuo
% ═══════════════════════════════════════════════════════════════

\begin{grammatikbox}[G1: El presente continuo \de{(Verlaufsform der Gegenwart)}]

\begin{multicols}{2}

\textbf{Bildung:} \es{estar} (konjugiert) + \es{gerundio}

\vspace{4pt}

\de{Wie im Englischen:} \textit{I \textbf{am reading}} = \es{Estoy leyendo}

\vspace{6pt}

\textbf{Gerundio-Endungen:}

\vspace{2pt}

\begin{tabular}{@{}ll@{}}
\es{-ar} → & \es{-ando} \de{(hablar → hablando)} \\
\es{-er/-ir} → & \es{-iendo} \de{(comer → comiendo)} \\
\end{tabular}

\columnbreak

\textbf{Estar:} \lsg{estoy}, estás, \lsg{está}, estamos, estáis, \lsg{están}

\vspace{6pt}

\textbf{Beispiel:}

\es{La chica \lsg{está} \lsg{leyendo}.}

\de{(Das Mädchen liest gerade.)}

\end{multicols}

\end{grammatikbox}

\vspace{4pt}

% ═══════════════════════════════════════════════════════════════
% AUFGABEN (zweispaltig)
% ═══════════════════════════════════════════════════════════════

\begin{multicols}{2}

% --- AUFGABE 1: Gerundio bilden ---
\aufgabe{A1}{Bilde das Gerundio.} \punktebox{6}

{\small Schreibe die Gerundio-Form der Verben.}

\vspace{4pt}

\begin{tabular}{@{}p{2.8cm}p{3.2cm}@{}}
\toprule
\textbf{Infinitivo} & \textbf{Gerundio} \\
\midrule
caminar \de{laufen} & \lsg{caminando} \\[4pt]
pasear \de{spazieren} & \lsg{paseando} \\[4pt]
mirar \de{schauen} & \lsg{mirando} \\[4pt]
hablar \de{sprechen} & \lsg{hablando} \\[4pt]
comer \de{essen} & \lsg{comiendo} \\[4pt]
beber \de{trinken} & \lsg{bebiendo} \\[4pt]
leer \de{lesen} & \lsg{leyendo} \\[4pt]
escribir \de{schreiben} & \lsg{escribiendo} \\[4pt]
jugar \de{spielen} & \lsg{jugando} \\[4pt]
correr \de{rennen} & \lsg{corriendo} \\[4pt]
reír \de{lachen} & \lsg{riendo} \\[4pt]
dormir \de{schlafen} & \lsg{durmiendo} \\
\bottomrule
\end{tabular}

\columnbreak

% --- AUFGABE 2: Bildteile zuordnen ---
\aufgabe{A2}{Ordne zu: Teile des Bildes.} \punktebox{4}

{\small Verbinde die deutschen und spanischen Begriffe.}

\vspace{6pt}

\begin{tikzpicture}[
    regelbox/.style={draw, rectangle, minimum width=2.4cm, minimum height=0.7cm,
                     font=\small, align=center},
    beispielbox/.style={draw, rectangle, minimum width=2.6cm, minimum height=0.7cm,
                        font=\small}
]
% Linke Spalte: Deutsch
\node[regelbox] (D1) at (0,0) {im Vordergrund};
\node[regelbox] (D2) at (0,-1) {im Hintergrund};
\node[regelbox] (D3) at (0,-2) {links};
\node[regelbox] (D4) at (0,-3) {in der Mitte};

% Rechte Spalte: Spanisch (gemischt)
\node[beispielbox] (S1) at (4.2,0) {a la izquierda};
\node[beispielbox] (S2) at (4.2,-1) {en el centro};
\node[beispielbox] (S3) at (4.2,-2) {en primer plano};
\node[beispielbox] (S4) at (4.2,-3) {al fondo};

% LÖSUNGEN: Verbindungslinien
\draw[loesung, thick] (D1.east) -- (S3.west);
\draw[loesung, thick] (D2.east) -- (S4.west);
\draw[loesung, thick] (D3.east) -- (S1.west);
\draw[loesung, thick] (D4.east) -- (S2.west);
\end{tikzpicture}

\vspace{8pt}

% --- AUFGABE 3: Positionen ---
\aufgabe{A3}{Ergänze die Positionen.} \punktebox{4}

{\small Wähle: \es{de pie · sentado/a · tumbado/a · apoyado/a}}

\vspace{4pt}

\begin{teil}
    \item El hombre está \lsg{sentado} en el banco.\\
          \de{Der Mann sitzt auf der Bank.}
    \item La mujer está \lsg{apoyada} en la pared.\\
          \de{Die Frau lehnt an der Wand.}
    \item Los niños están \lsg{tumbados} en el césped.\\
          \de{Die Kinder liegen auf dem Rasen.}
    \item Mi padre está \lsg{de pie} delante de la tienda.\\
          \de{Mein Vater steht vor dem Geschäft.}
\end{teil}

\end{multicols}

% ═══════════════════════════════════════════════════════════════
% SEITE 2
% ═══════════════════════════════════════════════════════════════

\newpage

% --- Header Seite 2 ---
\noindent
\begin{tikzpicture}[remember picture, overlay]
    \fill[kopfzeile] (current page.north west) rectangle +(\paperwidth, -0.7cm);
    \node[anchor=west, text=white, font=\sanstext\small\bfseries]
        at ([xshift=1.4cm, yshift=-0.35cm]current page.north west) {\thema{} – \textcolor{loesung}{MUSTERLÖSUNG} (Fortsetzung)};
\end{tikzpicture}

\vspace{0.4cm}

% --- AUFGABE 4: Musterbeschreibung analysieren ---
\aufgabe{A4}{Analysiere die Musterbeschreibung.} \punktebox{4}

{\small Markiere im Text: \textcolor{akzent}{rot} = Bildaufbau, \underline{unterstreichen} = presente continuo, \fbox{Kasten} = Position}

\vspace{6pt}

\begin{vokabelbox}[Musterbeschreibung]
\es{\textcolor{akzent}{En la foto} veo un parque en una ciudad. \textcolor{akzent}{En primer plano} hay un banco donde una mujer \fbox{está sentada}. Ella \underline{está leyendo} un libro. \textcolor{akzent}{A la izquierda}, un hombre \underline{está paseando} con su perro. \textcolor{akzent}{Al fondo} se ven varios edificios altos. \textcolor{akzent}{En el centro} de la imagen hay una fuente bonita. \textcolor{akzent}{A la derecha}, dos niños \underline{están jugando} en el césped.}
\end{vokabelbox}

\vspace{6pt}

\begin{multicols}{2}

% --- AUFGABE 5: Eigene Bildbeschreibung ---
\aufgabe{A5}{Schreibe eine Bildbeschreibung.} \punktebox{8}

{\small Beschreibe das Foto, das dein Lehrer dir gibt. Nutze:}

\begin{itemize}[leftmargin=1.2em, itemsep=1pt]
    \item Bildaufbau: \es{en primer plano, al fondo, a la izquierda...}
    \item Positionen: \es{estar de pie, sentado, tumbado...}
    \item Presente continuo: \es{está + gerundio}
    \item Personenbeschreibung: \es{alto, pelo rubio, ojos azules...}
\end{itemize}

\vspace{6pt}

\textcolor{loesung}{\textit{Individuelle Lösung. Bewertungskriterien:}}

\vspace{4pt}

\textcolor{loesung}{\footnotesize
\begin{itemize}[leftmargin=1.2em, itemsep=1pt]
    \item 2 P: Bildaufbau (mind. 2 Ausdrücke)
    \item 2 P: Presente continuo (mind. 2 korrekt)
    \item 1 P: Positionen (mind. 1 korrekt)
    \item 1 P: Personenbeschreibung vorhanden
    \item 1 P: Orte/Gegenstände benannt
    \item 1 P: Redemittel verwendet
\end{itemize}}

\columnbreak

% --- VOKABEL-HILFE (kompakt) ---
\begin{vokabelbox}[Vokabelhilfe: Orte]
{\small
\begin{tabular}{@{}ll@{}}
el parque & \de{der Park} \\
el banco & \de{die Bank} \\
el árbol & \de{der Baum} \\
el césped & \de{der Rasen} \\
la fuente & \de{der Brunnen} \\
el camino & \de{der Weg} \\
\midrule
la calle & \de{die Straße} \\
la plaza & \de{der Platz} \\
el edificio & \de{das Gebäude} \\
la tienda & \de{das Geschäft} \\
\midrule
la cafetería & \de{das Café} \\
la terraza & \de{die Terrasse} \\
la mesa & \de{der Tisch} \\
la silla & \de{der Stuhl} \\
\end{tabular}
}
\end{vokabelbox}

\vspace{6pt}

\begin{vokabelbox}[Redemittel]
{\small
\begin{tabular}{@{}l@{}}
\es{En la foto veo...} \\
\es{En primer plano hay...} \\
\es{Al fondo se ve...} \\
\es{Una persona está + gerundio...} \\
\es{Parece que...} \de{(Es scheint, dass...)} \\
\end{tabular}
}
\end{vokabelbox}

\end{multicols}

% ═══════════════════════════════════════════════════════════════
% K1: CHECKLISTE
% ═══════════════════════════════════════════════════════════════

\vspace{4pt}

\begin{checklistbox}
\sanstext{\small\textbf{K1: Checkliste – Gute Bildbeschreibung}}

\vspace{4pt}

{\small
\begin{multicols}{2}
$\boxtimes$ Bildaufbau verwendet (primer plano, al fondo...)\\
$\boxtimes$ Presente continuo korrekt (estar + gerundio)\\
$\boxtimes$ Positionen beschrieben (de pie, sentado...)
\columnbreak
$\boxtimes$ Personen beschrieben (pelo, ojos, alto...)\\
$\boxtimes$ Orte/Gegenstände benannt\\
$\boxtimes$ Redemittel genutzt (En la foto veo...)
\end{multicols}
}
\end{checklistbox}

% ═══════════════════════════════════════════════════════════════
% FOOTER
% ═══════════════════════════════════════════════════════════════

\vfill

\noindent\rule{\textwidth}{0.5pt}

\vspace{4pt}

\hfill\sanstext{\textbf{Punkte:}} \luecke{1.2cm} / \maxpunkte

\end{document}
